% Quy trình sinh câu trả lời có tra cứu tài liệu
\begin{tikzpicture}[baocao]
  \node[hop, text width=2.0cm] (q) at (1.2,0) {Câu hỏi của người dùng};
  \node[hopxam, text width=2.3cm] (r) at (4.5,0) {\textbf{Bước 1}\\Tìm các đoạn tài liệu liên quan};
  \node[hopxam, text width=2.3cm] (g) at (7.9,0) {\textbf{Bước 2}\\Đưa các đoạn tìm được vào lời nhắc};
  \node[hopxam, text width=2.3cm] (l) at (11.3,0) {\textbf{Bước 3}\\Mô hình viết câu trả lời};
  \node[hop, text width=1.9cm] (a) at (14.4,0) {Câu trả lời kèm nguồn};
  \node[hop, text width=2.3cm] (k) at (4.5,-2.9) {Kho tài liệu đã lập chỉ mục};
  \draw[mui] (q) -- (r);
  \draw[mui] (r) -- (g);
  \draw[mui] (g) -- (l);
  \draw[mui] (l) -- (a);
  \draw[haichieu] (k) -- (r);
\end{tikzpicture}
